% !TEX root = ../notes.tex

\documentclass[../notes.tex]{subfiles}

\begin{document}

\section{April 6}
Here we go. The quiz on Wednesday will have two problems, and we will be able to use the entire class period.

\subsection{Cusp Forms}
Here is our definition.
\begin{definition}[cuspidal]
	Fix a reductive group $G$ over a number field $F$. Then an automorphic form $f$ is \textit{cuspidal} if and only if
	\[\int_{N(F)\backslash N(\AA_F)}f(ng)\,dn=0\]
	for any $g\in G(\AA_F)$ and proper parabolic subgroup $P\subseteq G$ (defined over $F$) with unipotent radical $N$. We let $\mc A_{\mathrm{cusp}}$ denote the subspace of cusp forms.
\end{definition}
\begin{remark}
	The integral converges because $f$ is smooth, and the space $N(F)\backslash N(\AA_F)$ is compact. Indeed, $N$ is an iterated extension of $\mathbb G_a$s, so $N(F)\backslash N(\AA_F)$ is an iterated extension of the compact spaces $F\backslash\AA_F$.
\end{remark}
\begin{remark}
	It is enough to only consider the maximal parabolic subgroups: a larger parabolic subgroup produces a smaller unipotent radical, so the integral for the larger parabolic can be tiled by integrals for the larger parabolic. It is also enough to even check a single representative in each conjugacy class. For $\op{GL}_n$, these look like the block-upper triangular matrices
	\[\begin{bmatrix}
		* & * \\ & *
	\end{bmatrix}.\]
	A non-maximal such parabolic subgroup may have more blocks.
\end{remark}
\begin{exe}
	Choose a modular form $f$ of weight $k$ and level $\op{SL}_2(\ZZ)$, and let
	\[\varphi_f\left(\begin{bmatrix}
		a & b \\ c  & d
	\end{bmatrix}\right)\coloneqq f\left(\frac{ai+b}{ci+d}\right)(ci+d)^{-k}\]
	be the corresponding automorphic form. We will check that $f$ is cuspidal if and only if $\varphi_f$ is.
\end{exe}
\begin{proof}
	The only parabolic to check (up to conjugacy) is the Borel subgroup, so we let $N$ be the subgroup $\begin{bsmallmatrix}
		1 & * \\ & 1
	\end{bsmallmatrix}$, and we would like to check if the integral
	\[\int_{N(\QQ)\backslash N(\AA_\QQ)}\varphi_f(ng)\,dn=\int_{\ZZ\backslash\RR}f(z+x)\,dx\]
	vanishes, where $z=gi$. Namely, one can see that the integral on the left descends to $\QQ\backslash\AA_\QQ/\widehat{\ZZ}=\ZZ\backslash\RR$, which is the right-hand integral. Writing $f(z)=\sum_na_nq^n$ with $q=e^{2\pi iz}$, we thus see that this integral computes the constant term $a_0$ in this Fourier expansion: $a_0=\frac1{2\pi i}\oint_{\gamma}f(q)\,\frac{dq}q$, where $\gamma$ is the corresponding loop around $0$. Thus, we see that $\varphi_f$ is cuspidal if and only if $a_0=0$, which is what it means for $f_0$ to be cuspidal.
\end{proof}
Cusp forms are important for us because they have an application to the spectral decomposition. In particular, we are interested in decomposing $L^2(G(\QQ)\backslash G(\AA_\QQ))$.
\begin{lemma}
	Fix a semisimple group $G$ over a number field $F$. Then the space of automorphic forms embeds as a dense subspace of $L^2(G(\QQ)\backslash G(\AA_\QQ))$.
\end{lemma}
\begin{proof}
	We will not show this. Note that the embedding exists because $G(\QQ)\backslash G(\AA_\QQ)$ has finite volume.
\end{proof}
\begin{remark}
	A similar statement works for $G$ reductive. There is a slight caveat: 
\end{remark}
\begin{notation}
	Fix a semisimple group $G$ over $\QQ$. Then we let $L^2_{\mathrm{cusp}}(G(\QQ)\backslash G(\AA_\QQ))$ be the completion of $\mc A_{\mathrm{cusp}}$ in $L^2(G(\QQ)\backslash G(\AA_\QQ))$.
\end{notation}
Now, $G(\AA_\QQ)$ acts by right translation on $L^2(G(\QQ)\backslash G(\AA_\QQ))$, and it preserves the integration pairing given by the Haar measure. Thus, we are basically looking at some operators on a Hilbert space, where one expects to have a discrete spectrum and a continuous spectrum. This cuspidal subspace turns out to give us a lot of information about the discrete spectrum.
\begin{proposition}
	Fix a semisimple group $G$ over $\QQ$. Then the action of $G(\AA_\QQ)$ on $L^2_{\mathrm{cusp}}(G(\QQ)\backslash G(\AA_\QQ))$ decomposes into a direct sum of irreducible representations with finite multiplicity.
\end{proposition}
\begin{example}
	For $G=\op{GL}_n$, irreducible representations appear at most once.
\end{example}
\begin{remark}
	There may be other subspaces in the ``discrete spectrum.'' For example, the constant functions are certainly in the discrete spectrum. For $G=\op{SL}_2$, it turns out that the discrete spectrum is the sum of the cuspidal part and the constant functions.
\end{remark}
We now turn to the continuous spectrum. Here, these are constructed by parabolic induction, using Eisenstein series. The idea, roughly speaking, is to take a parabolic subgroup $P\subseteq G$ with Levi subgroup $M$. Then choose a discrete irreducible automorphic representation $\pi$ of $M$, which embeds into $L^2(M(\QQ)\backslash M(\AA_\QQ);\chi)$, where $\chi$ is some fixed central character. From here, there is a process which allows us to induce from $L^2(M(\QQ)\backslash M(\AA_\QQ);\chi)$ back to $G$.
\begin{example}
	Suppose that $P$ is a Borel subgroup so that $M$ is a torus $T$. Then the choice of central character amounts to a choice of character of $T$. If $T$ is split of rank $n$, then we are basically hunting for $n$ id\'ele class characters, so there are $n$ continuos parameters in this ``continuous spectrum.''
\end{example}
% For example, if $G$ is semisimple over $\QQ$, then it turns out that $\mc A$ embeds a dense subspace of $L^2(G(\QQ)\backslash G(\AA_\QQ))$, and we let $L^2_{\mathrm{cusp}}(G(\QQ)\backslash G(\AA_\QQ))$ be the completion of $\mc A_{\mathrm{cusp}}$. Now, $G(\AA_\QQ)$ acts by right translation on $L^2(G(\QQ)\backslash G(\AA_\QQ))$, and it preserves the integration pairing given by the Haar measure. It turns out that the action of $G(\AA_\QQ)$ on $L^2_{\mathrm{cusp}}(G(\QQ)\backslash G(\AA_\QQ))$ decomposes into a direct sum of irreducible representations with finite multiplicity (though there is a continuous spectrum for the leftover pieces in general). For example, f

\subsection{Automorphic Forms for Function Fields}
Let's work with function fields. For convenience, let $G$ be a split semisimple group over a finite field $\FF_q$, and we let $X$ be a smooth projective geometrically integral curve, and we set $F\coloneqq\FF_q(X)$.\footnote{Because we are assuming that $G$ is split, it already has a model over $\ZZ$. Thus, we do not lose anything by descending $G$ to $\FF_q$ instead of over $X$.} Our automorphic forms should approximately come from $C^\infty(G(F)\backslash G(\AA_F))$, which already has an action by $G(\AA_F)$ by translation.

Here is our definition of automorphic form, which notably does not have to fight with analytic difficulties to get off the ground.
\begin{definition}[automorphic form]
	Fix a function field $F$, and let $G$ be a split semisimple group. Let $G(\AA_F)$ act on $C^\infty(G(F)\backslash G(\AA_F))$ by right translations, denoted $R$. Then $f\in C^\infty(G(F)\backslash G(\AA_F))$ is an \textit{automorphic form} if and only if
	\[\op{span}\{R(g)f:g\in G(\AA_F)\}\]
	is an admissible representation of $G(\AA_F)$.
\end{definition}
Unsurprisingly, we have a notion of cusp forms.
\begin{definition}[cuspidal]
	Fix a function field $F$, and let $G$ be a split semisimple group. An automorphic form $f$ is \textit{cuspidal} if and only if
	\[\int_{N(F)\backslash N(\AA_F)}f(ng)\,dn=0\]
	for all $g\in G(\AA_F)$ and all proper parabolic subgroups $P\subseteq G$, whose unipotent radical is denoted $N$.
\end{definition}
The benefit of working over function fields is that we can do geometry. To understand automorphic forms, we may as well start by fixing an invariant compact open subgroup. For example, for a divisor $D=\sum_xn_x[x]$, we may take
\[K_D\coloneqq\prod_{x\in\left|X\right|}\Gamma_x(n_x),\]
where $\Gamma_x(n_x)$ is the kernel of the natural projection $G(\OO_x)\to G(\OO_x/\mf m_x^{n_x})$. Thus, we are now interested in studying functions on the space $G(F)\backslash G(\AA_F)/K_D$.

On the homework, we have already seen a version of this interpretation for number fields.
\begin{example}
	For $\op{GL}_n$ over $\QQ$, one can interpret $\op{GL}_n(\QQ)\backslash\op{GL}_n(\AA_F)/(\op O_n(\QQ)\times\op{GL}_n(\widehat\ZZ))$ as a pair $(\Lambda,q)$ of a projective $\ZZ$-module $\Lambda$ of rank $n$ equipped with a positive-definite quadratic form $q$ of $\Lambda_\RR$.
\end{example}
One can do the same thing over function fields.
\begin{exe}
	Fix a function field $F=\FF_q(X)$. Then $\op{GL}_n(F)\backslash\op{GL}_n(\AA_F)/K_0$ is equivalent (as a groupoid) to the category of vector bundles of rank $n$ on $X$.
\end{exe}
\begin{proof}
	It is enough to show $\op{GL}_n(\AA_F)/K_0$ is equivalent to the category of pairs $(\mc E,\tau)$ where $\mc E$ is a vector bundle of rank $n$ on $X$, and $\tau$ is a trivialization $\mc E|_{\Spec F}\cong F^n$ (at the generic fiber). (Here, $\op{GL}_n(F)$ acts on the target category by acting on the target $F^n$ of $\tau$.) To define our map, we note that $\op{GL}_n(\AA_F)/K_0$ is
	\[\bigoplus_x\op{GL}_n(F_x)/\op{GL}_n(\OO_x),\]
	and we know that $\op{GL}_n(F_x)/\op{GL}_n(\OO_x)$ classifies $\OO_x$-lattices embedded in $F_x^n$. Thus, we may define our map by sending $(\mc E,\tau)$ in the target to the lattices $\tau(\mc E_{\OO_x})\subseteq F_x^n$. Here are our checks on this map.
	\begin{itemize}
		\item A determinant calculation shows that $\tau$ spreads out as an isomorphism to $\OO_x^n$ to all but finitely many places (namely, the determinant only needs finitely many denominators). Thus, $\tau(\mc E_{\OO_x})\subseteq F_x^n$ is the standard lattice $\OO_x^n$ for all but finitely many $x$.
		\item We will only sketch the remaining checks.  To go backward, one needs to be able to construct a vector bundle from the data of our local lattices. The point is to start from 
		\qedhere
	\end{itemize}
\end{proof}
\begin{remark}
	The same argument can show that $\op{GL}_n(F)\backslash\op{GL}_n(\AA_F)/K_D$ is equivalent (as a groupoid) to the category of vector bundles $\mc E$ of rank $n$ along with a level structure isomorphism $\tau\colon\mc E|_{\Spec\OO_D}\cong\OO_D^n$.
\end{remark}
\begin{remark}
	One can give similar interpretations for other classical groups. For example, for $G=\op{Sp}_{2n}$, one replaces vector bundles with ``symplectic vector bundles,'' which are vector bundles equipped with a symplectic pairing.
\end{remark}
The moral is that automorphic forms can be interpreted as functions on a space of ``$G$-bundles.''

While we're here, let's also give a geometric interpretation of cusp forms of level $K_0$. Thus, we would like a geometric interpretation of the integral
\[\int_{N(F)\backslash N(\AA_F)}f(ng)\,dn,\]
where $N$ is the unipotent radical of some proper parabolic subgroup $P$.
\begin{example}
	Take $G=\op{GL}_2$. Then $B(F)\backslash B(\AA_F)/(K_0\cap B(\AA_F))$ embeds into $G(F)\backslash G(\AA_F)/K_0$, and one can show that the former classifies vector bundles $\mc E$ of rank $2$ along with a vector bundle $\mc L\subseteq\mc E$ of rank $1$ for which $\mc E/\mc L$ is also a vector bundle. Now, let $B=TN$ be the Levi decomposition, and we see that $N(F)\backslash N(\AA_F)/(K_0\cap N(\AA_F))$ can be seen as the fiber of the map
	\[B(F)\backslash B(\AA_F)/(K_0\cap N(\AA_F))\to T(F)\backslash T(\AA_F)/(K_0\cap T(\AA_F)).\]
	It thus turns out that $N(F)\backslash N(\AA_F)/(K_0\cap N(\AA_F))$ parameterizes pairs short exact sequences $0\to\OO_X\to\mc E\to\OO_X\to0$. Indeed, we are asking for the pair $(\mc E,\mc L)$ from $\mc E$ to have a given trivialization of both $\mc L$ and $\mc E/\mc L$.
\end{example}
\begin{example}
	We continue with  $G=\op{GL}_2$. Then, we see that $\int_{N(F)\backslash N(\AA_F)}f(n)\,dn$ is a suitable average of $f(\mc E)$, where $\mc E$ varies over short exact sequences
	\[0\to\OO_X\to\mc E\to\OO_X\to0.\]
	Notably, a single $\mc E$ (up to isomorphism of vector bundles) may appear multiple times in the sum because it may be in different isomorphism classes of short exact sequence. In general, for $\int_{N(F)\backslash N(\AA_F)}f(ng)$, we may use the Iwasawa decomposition $G(\AA_F)=N(\AA_F)T(\AA_F)K_0$ to force $g$ to be in $T(\AA_F)$. In fact, we even only care about $g$ as a class in $T(F)\backslash T(\AA_F)/(K_0\cap T(\AA_F))=\op{Pic}(X)^2$, which means that $g$ is represented by a pair $(\mc L_1,\mc L_2)$ of line bundles. It follows that $f$ is a cusp form if and only if we always have
	\[\sum_{\mc E\in\op{Ext}^1(\mc L_1,\mc L_2)}f(\mc E)=0.\]
\end{example}
\begin{remark}
	For $G=\op{GL}_n$, a similar argument shows that $f$ is a cusp form if and only if we have
	\[\sum_{\mc E\in\op{Ext}^1(\mc E_1,\mc E_2)}f(\mc E)=0\]
	for all pairs $(\mc E_1,\mc E_2)$ of vector bundles of rank summing to $n$.
\end{remark}
Next time, we will sketch the following theorem for $\op{GL}_n$.
\begin{theorem}
	Fix a split semisimple group $G$ over a function field $F$, and let $K\subseteq G(\AA_F)$ be a compact open subgroup. Then $C_{\mathrm{cusp}}(G(F)\backslash G(\AA_F)/K)$ is finite-dimensional.
\end{theorem}
This is some version of the admissibility for this cuspidal subspace.

\end{document}